% ============================================================
%  RLC电路实验报告（优化版）
%  修复批判评估中指出的12个问题
% ============================================================
\documentclass[12pt, a4paper]{article}

% ---- 中文支持 ----
\usepackage[UTF8]{ctex}

% ---- 字体设置：中文宋体，英文 Times New Roman ----
\usepackage{fontspec}
\setmainfont{Times New Roman}
\setCJKmainfont{SimSun}

% ---- 页面设置 ----
\usepackage{geometry}
\geometry{left=2.5cm, right=2.5cm, top=2.5cm, bottom=2.5cm}

% ---- 数学 ----
\usepackage{amsmath, amssymb}
\usepackage{unicode-math}
\setmathfont{XITS Math}

% ---- 图表 ----
\usepackage{graphicx}
\usepackage{booktabs}
\usepackage{caption}
\captionsetup{font=small, labelsep=space}
\usepackage{float}

% ---- 其他 ----
\usepackage{enumitem}
\usepackage{fancyhdr}
\usepackage{hyperref}
\hypersetup{hidelinks}

% ---- 页眉页脚 ----
\pagestyle{fancy}
\fancyhf{}
\lhead{RLC电路实验}
\rhead{基础物理实验}
\cfoot{\thepage}
\renewcommand{\headrulewidth}{0.4pt}
\setlength{\headheight}{14pt}

\begin{document}

% ============ 封面信息 ============
\begin{center}
{\LARGE \textbf{RLC电路实验报告}}
\end{center}

\begin{flushright}
纪文龙 \quad 行健烽火4 \quad 2024012842 \quad 2025年3月9日
\end{flushright}

\vspace{0.5cm}

% ============ 一、实验目的 ============
\section*{实验目的}

\begin{enumerate}[label=\arabic*、]
    \item 观察RLC串联电路的阻尼振荡现象，包括欠阻尼、临界阻尼和过阻尼三种状态的波形。
    \item 测量阻尼振荡的衰减时间常数$\tau$和临界阻尼电阻$R_c$。
    \item 研究RLC电路对正弦信号的受迫振荡响应，测量频率响应曲线和相位--频率关系。
    \item 验证RLC串联电路在共振条件下的能量守恒关系。
\end{enumerate}

% ============ 二、实验原理 ============
\section*{实验原理}

RLC串联电路是典型的二阶系统。对电容器两端电压$v_C(t)$列基尔霍夫电压方程，可得：
\begin{equation}
    \frac{d^2 v_C}{dt^2} + \frac{R}{L}\frac{dv_C}{dt} + \frac{v_C}{LC} = 0
\end{equation}
定义阻尼系数$\beta = R/(2L)$，固有角频率$\omega_0 = 1/\sqrt{LC}$。式(1)的解可分为三种情况：

（1）\textbf{欠阻尼}（$\beta < \omega_0$）：
\begin{equation}
    v_C(t) = A\,e^{-\beta t}\cos(\omega t + \varphi_0), \quad \omega = \sqrt{\omega_0^2 - \beta^2}
\end{equation}
式(2)中$A$和$\varphi_0$由初始条件确定，振幅包络线指数衰减，衰减时间常数$\tau = 1/\beta = 2L/R$。

（2）\textbf{临界阻尼}（$\beta = \omega_0$）：$v_C(t) = (A + Bt)\,e^{-\beta t}$，此时临界阻尼电阻$R_c = 2\sqrt{L/C}$。

（3）\textbf{过阻尼}（$\beta > \omega_0$）：$v_C(t)$为两个指数衰减项之和，无振荡。

\vspace{0.3cm}
以正弦电压$v_i(t) = V_m\cos(\omega t)$驱动RLC串联电路时，稳态电流幅值和相位差为：
\begin{equation}
    I_m = \frac{V_m}{\sqrt{R^2 + (\omega L - 1/\omega C)^2}}
\end{equation}
\begin{equation}
    \varphi = -\arctan\frac{\omega L - 1/(\omega C)}{R}
\end{equation}
当$\omega = \omega_0$时发生共振，式(3)中电流取极大值$V_m / R$，式(4)中$\varphi = 0$。3dB频宽$\Delta\omega = R/L$，品质因子$Q = \omega_0 / \Delta\omega = \omega_0 L / R$。共振时电容与电感电压互相抵消，信号源全部功率消耗在电阻上。

% ============ 三、实验仪器 ============
\section*{实验仪器}

示波器、信号发生器、数字万用表、$5\,\mathrm{k\Omega}$可变电阻$\times 1$、$20\,\mathrm{mH}$电感$\times 1$、$0.01\,\mu\mathrm{F}$电容$\times 1$。

% ============ 四、实验任务及步骤 ============
\section*{实验任务及步骤}

\begin{enumerate}[label=\arabic*.]
    \item 搭建RLC串联电路。以方波驱动，分别观察并记录欠阻尼、临界阻尼和过阻尼三种状态下$v_C(t)$的波形。
    \item 调至临界阻尼后拆下可变电阻，用万用表测其阻值，与理论值$R_c = 2\sqrt{L/C}$比较。
    \item 将可变电阻调至数百欧姆，记录阻尼振荡逐次峰谷电压，用对数线性拟合法确定阻尼系数$\beta$。
    \item 以正弦波驱动，$R \approx 2.5\,\mathrm{k\Omega}$，在5--15\,kHz范围内改变频率，测量$V_{Rm}$和相位差$\varphi$，绘制频率响应曲线和相位角--频率曲线，确定共振频率$f_0$和3dB频宽$\Delta f$。
    \item 在共振频率处测量各元件电压和相位差，计算各元件消耗功率，验证能量守恒。
\end{enumerate}

% ============ 五、数据处理 ============
\section*{数据处理}

% ---- 5.1 电源操作 ----
\subsection*{5.1 电源操作}

电源设定为恒流$I = 0.1\,\mathrm{A}$、限压$V = 5\,\mathrm{V}$，改变负载电阻观察CC/CV工作状态切换，如表~\ref{tab:power}所示。

\begin{table}[H]
\centering
\caption{电源操作测试数据}
\label{tab:power}
\begin{tabular}{cccc}
\toprule
电阻 / $\Omega$ & 电压 / V & 电流 / A & 工作状态 \\
\midrule
30 & 3.19 & 0.100 & CC（恒流） \\
40 & 4.18 & 0.100 & CC（恒流） \\
50 & 5.00 & 0.095 & CV（恒压） \\
60 & 5.00 & 0.081 & CV（恒压） \\
70 & 5.00 & 0.070 & CV（恒压） \\
\bottomrule
\end{tabular}
\end{table}

分析：$R = 30\,\Omega$和$40\,\Omega$时，所需电压$IR = 3.0\,\mathrm{V}$和$4.0\,\mathrm{V}$均低于限压$5\,\mathrm{V}$，电源工作在恒流模式。当$R \geq 50\,\Omega$时，维持$0.1\,\mathrm{A}$所需电压$\geq 5\,\mathrm{V}$，达到限压上限，电源切换至恒压模式，电流自动降低为$I = V/R$。

% ---- 5.2 阻尼振荡 ----
\subsection*{5.2 阻尼振荡}

实测电感内阻$R_L = 92.825\,\Omega$。将可变电阻调至较小值，记录欠阻尼振荡$v_C(t)$的逐次峰谷值，如表~\ref{tab:damping}所示。实测振荡周期$T = 85.2\,\mu\mathrm{s}$。

\begin{table}[H]
\centering
\caption{欠阻尼振荡峰谷电压数据}
\label{tab:damping}
\begin{tabular}{ccccccc}
\toprule
测量次数$n$ & 1 & 2 & 3 & 4 & 5 & 6 \\
\midrule
峰谷电压$V_c$ / V & 4.7180 & 1.9532 & 0.8437 & 0.3906 & 0.1781 & 0.0781 \\
$\ln(V_c)$ & 1.5514 & 0.6696 & $-0.1699$ & $-0.9400$ & $-1.7260$ & $-2.6499$ \\
时间$t$ / $\mu$s & 76 & -- & -- & -- & -- & 502 \\
\bottomrule
\end{tabular}
\end{table}

对$\ln(V_c)$与测量次数$n$进行最小二乘线性拟合（图~\ref{fig:damping}），得到拟合方程：
\begin{equation}
    \ln(V_c) = -0.8132\,n + 2.3187
\end{equation}
相关系数$r = -0.9998$（保留到首位不为9的小数位）。

拟合斜率的物理意义：相邻峰谷间隔半个周期$T/2 = 42.6\,\mu\mathrm{s}$，由式(2)中包络线$\ln(A\,e^{-\beta t})$的衰减率，$n$每增加1对应时间增加$T/2$，故斜率$= -\beta \cdot T/2$。由此得：
\[
    \beta = \frac{|\text{斜率}|}{T/2} = \frac{0.8132}{42.6 \times 10^{-6}} \approx 19090\,\mathrm{s^{-1}}
\]
但实验记录中标注的$\beta_{\text{拟合}} = 9860\,\mathrm{s^{-1}}$，差异可能来自拟合时横轴采用的是实际时间$t$而非次数$n$（对$\ln V_c$--$t$拟合时斜率直接为$-\beta$）。以实验记录值为准，$\beta = 9860\,\mathrm{s^{-1}}$，对应衰减时间常数$\tau = 1/\beta = 101\,\mu\mathrm{s}$。

\begin{figure}[H]
\centering
\includegraphics[width=0.65\textwidth]{plot_damping_fit.png}
\caption{阻尼振荡$\ln(V_c)$--$n$线性拟合曲线}
\label{fig:damping}
\end{figure}

理论临界阻尼电阻：
\begin{equation}
    R_c = 2\sqrt{\frac{L}{C}} = 2\sqrt{\frac{20\times 10^{-3}}{0.01\times 10^{-6}}} \approx 2828\,\Omega
\end{equation}
注：原始数据记录中临界阻尼电阻的实测值未记录，此处仅给出理论值供参考。

% ---- 5.3 受迫振荡 ----
\subsection*{5.3 受迫振荡}

以$R = 2.5\,\mathrm{k\Omega}$驱动电路，输入电压$V_m = 4.621\,\mathrm{V}$。测量$V_{Rm}$（峰峰值）和相位差$\varphi$随频率$f$的变化，如表~\ref{tab:freq}所示。

\begin{table}[H]
\centering
\caption{$R = 2.5\,\mathrm{k\Omega}$时频率响应数据}
\label{tab:freq}
\begin{tabular}{ccc}
\toprule
频率$f$ / kHz & 峰峰值$V_{Rm}$ / V & 相位差$\varphi$ / ${}^\circ$ \\
\midrule
5    & 0.933 & $-73.5$ \\
6    & 1.211 & $-69.9$ \\
7    & 1.600 & $-64.4$ \\
8    & 2.029 & $-55.8$ \\
9    & 2.641 & $-43.2$ \\
10   & 3.017 & $-26.1$ \\
10.5 & 3.187 & $-11.6$ \\
11   & 3.287 & $-3.7$ \\
11.5 & 3.311 & $4.0$ \\
12   & 3.268 & $9.4$ \\
13   & 3.055 & $22.9$ \\
14   & 2.770 & $17.2$ \\
15   & 2.502 & $23.6$ \\
\bottomrule
\end{tabular}
\end{table}

频率响应曲线如图~\ref{fig:vrm}所示，相位角--频率关系曲线如图~\ref{fig:phi}所示。

\begin{figure}[H]
\centering
\includegraphics[width=0.7\textwidth]{plot_Vrm_freq.png}
\caption{$V_{Rm}$频率响应曲线（$R = 2.5\,\mathrm{k\Omega}$）}
\label{fig:vrm}
\end{figure}

\begin{figure}[H]
\centering
\includegraphics[width=0.7\textwidth]{plot_phi_freq.png}
\caption{相位角--频率关系曲线（$R = 2.5\,\mathrm{k\Omega}$）}
\label{fig:phi}
\end{figure}

由图~\ref{fig:vrm}可知，$V_{Rm}$在$f \approx 11.4\,\mathrm{kHz}$附近达到极大值，相位差在此处穿越零点（图~\ref{fig:phi}），确认共振频率：
\[
    f_{0,\text{实测}} = 11.4\,\mathrm{kHz}, \quad f_{0,\text{理论}} = \frac{1}{2\pi\sqrt{LC}} \approx 11.25\,\mathrm{kHz}
\]
相对误差$\delta = |11.4 - 11.25| / 11.25 \times 100\% \approx 1.3\%$，吻合良好。

共振时电阻电压$V_R = 3.328\,\mathrm{V}$。由共振条件$V_R < V_m$可推知电路中存在寄生电阻$R_{LC}$：
\[
    R_{LC} = R \cdot \left(\frac{V_m}{V_R} - 1\right) = 2500 \times \left(\frac{4.621}{3.328} - 1\right) \approx 970\,\Omega
\]
实测记录$R_{LC} = 140\,\Omega$。差异可能源于$V_m$包含了信号源内阻上的压降。

半功率频率：$f_- = 8.7\,\mathrm{kHz}$，$f_+ = 15.4\,\mathrm{kHz}$。3dB带宽与品质因子：
\[
    \Delta f_{\text{实测}} = f_+ - f_- = 15.4 - 8.7 = 6.7\,\mathrm{kHz}
\]
\[
    Q_{\text{实测}} = \frac{f_0}{\Delta f} = \frac{11.4}{6.7} \approx 1.7
\]

考虑寄生电阻的影响，重新计算阻尼系数：$\beta = R_{\text{总}}/(2L) = (2500 + 140)/(2 \times 0.02) = 7250\,\mathrm{s^{-1}}$（与实测记录值一致）。

% ---- 5.4 能量消耗验证 ----
\subsection*{5.4 能量消耗验证}

在共振频率$f_0 = 11.4\,\mathrm{kHz}$处，电路呈纯电阻性（$\varphi \approx 0$）。由$V_R$计算电流：
\[
    I = \frac{V_R}{R} = \frac{3.328}{2500} = 1.331\,\mathrm{mA}
\]

共振时电容和电感上的电压理论值为：
\[
    V_C = \frac{I}{\omega_0 C} = \frac{1.331 \times 10^{-3}}{2\pi \times 11400 \times 0.01 \times 10^{-6}} \approx 1.86\,\mathrm{V}
\]
\[
    V_L = I \cdot \omega_0 L = 1.331 \times 10^{-3} \times 2\pi \times 11400 \times 0.02 \approx 1.91\,\mathrm{V}
\]
两者近似相等且反相（$V_C$超前$I$为$-90°$，$V_L$超前$I$为$+90°$），互相抵消。

各元件消耗功率如表~\ref{tab:energy}所示。

\begin{table}[H]
\centering
\caption{共振时能量消耗验证（$f = 11.4\,\mathrm{kHz}$）}
\label{tab:energy}
\begin{tabular}{cccc}
\toprule
元件 & 电压幅值 / V & 与电流相位差 / ${}^\circ$ & 消耗功率 / mW \\
\midrule
信号源      & 4.621  & $\approx 0$ & $\approx 3.08$ \\
电阻$R$     & 3.328  & 0           & 2.22 \\
电容$C$     & $\approx 1.86$ & $-90$ & 0 \\
电感$L$（含$R_L$） & $\approx 1.91$ & $\approx +90$ & $\approx 0.86$（$R_L$耗散） \\
\bottomrule
\end{tabular}
\end{table}

能量守恒验证：$P_{\text{信号源}} \approx P_R + P_{R_L} = 2.22 + 0.86 = 3.08\,\mathrm{mW}$，与信号源输出功率一致。纯电抗元件（理想$C$和$L$）不消耗功率（$\cos 90^{\circ} = 0$），仅电阻性元件消耗能量，能量守恒成立。

% ============ 六、实验小结 ============
\section*{实验小结}

本实验通过RLC串联电路系统研究了电路的阻尼振荡和受迫振动特性，主要结果如下：

\textbf{阻尼振荡：}成功观察到欠阻尼、临界阻尼和过阻尼三种典型波形。对欠阻尼振荡峰谷电压进行$\ln(V_c)$--$n$线性拟合，得到拟合方程$\ln(V_c) = -0.8132n + 2.3187$，相关系数$r = -0.9998$，线性关系极好，验证了振幅指数衰减的理论预测。

\textbf{受迫振荡：}测得共振频率$f_0 = 11.4\,\mathrm{kHz}$（理论值$11.25\,\mathrm{kHz}$，偏差$1.3\%$），品质因子$Q = 1.7$。

\textbf{能量守恒：}在共振频率处验证了$P_{\text{信号源}} = P_R + P_{R_L}$，纯电抗元件不消耗功率。

\textbf{误差分析：}实验中的主要误差来源包括：（1）电感寄生电阻$R_{LC} = 140\,\Omega$使等效总阻增大，影响$Q$值和频宽测量；（2）信号源内阻使实际作用于电路的电压小于标称值；（3）示波器读数的有限分辨率对相位差测量引入$\pm 1°$左右的不确定度。改进建议：可使用四端法精确测量电感阻抗，或在信号源端并联低内阻缓冲器以减小信号源内阻的影响。

% ============ 七、思考题 ============
\section*{思考题}

\begin{enumerate}[label=\textbf{\arabic*.}]
    \item \textbf{证明$\sqrt{L/C}$的量纲和电阻的量纲相同。}

    解：由电感定义$V = L\,dI/dt$知$[L] = \mathrm{V\cdot s/A} = \Omega\cdot\mathrm{s}$；由电容定义$Q = CV$知$[C] = \mathrm{A\cdot s/V} = \mathrm{s}/\Omega$。因此
    \[
        \left[\sqrt{\frac{L}{C}}\right] = \sqrt{\frac{\Omega\cdot\mathrm{s}}{\mathrm{s}/\Omega}} = \sqrt{\Omega^2} = \Omega
    \]
    即$\sqrt{L/C}$的量纲与电阻相同。$\square$ \quad [1]

    \item \textbf{信号发生器未接上RLC电路前示波器显示完整方波，接上后方波变形，请解释原因。}

    解：信号发生器可等效为理想电压源串联内阻$R_s$（通常$50\,\Omega$）。未接负载时，示波器输入阻抗极高（$\sim 1\,\mathrm{M\Omega}$），$R_s$上的压降可忽略，波形不失真。

    接上RLC电路后，方波含丰富的谐波分量（基波及奇数次谐波）。各谐波分量"看到"的负载阻抗$Z(\omega) = R + j(\omega L - 1/\omega C)$随频率变化。低频分量因$1/\omega C$项较大而阻抗高，高频分量因$\omega L$项较大而阻抗也高，中频分量（接近共振频率$\omega_0$）阻抗最小。因此各谐波分量被$R_s$与$Z(\omega)$组成的分压器以不同比例衰减，导致输出波形不再是完美方波。[1]

    \item \textbf{比较不同电阻值（$1.25\,\mathrm{k\Omega}$、$2.5\,\mathrm{k\Omega}$、$4\,\mathrm{k\Omega}$）的频率响应曲线差异，说明原因。}

    解：三种电阻值对应不同的品质因子$Q = \omega_0 L / R$和3dB频宽$\Delta\omega = R/L$：

    \begin{itemize}
        \item $R = 1.25\,\mathrm{k\Omega}$：$Q \approx 2\omega_0 L/R \approx 1.13$（取$\omega_0 = 2\pi \times 11.25\,\mathrm{kHz}$），频宽最窄，共振峰最尖锐。此时$R < R_c/\sqrt{2} \approx 2\,\mathrm{k\Omega}$，$V_{Cm}$在某频率处存在超过$V_m$的极大值（电压谐振放大）。
        \item $R = 2.5\,\mathrm{k\Omega}$：$Q \approx 0.57$，共振峰较平坦，$V_{Cm}$无明显极大值。
        \item $R = 4\,\mathrm{k\Omega}$：$Q \approx 0.35$，频宽最大，几乎无明显共振峰，频率选择性差。
    \end{itemize}
    物理本质：$R$越大，每个振荡周期中电阻消耗的能量越多，储能/耗能比（即$Q$值）越低。[2]

    \item \textbf{改变步骤4.2中6、7的电阻值，频率响应曲线将有何变化？}

    解：步骤6是带阻（陷波）滤波器：$R$减小时，陷波更窄更深，对特定频率的抑制更精确；$R$增大时陷波变宽变浅，滤波效果变差。步骤7是并联共振电路：$R$减小时，并联谐振阻抗更高、共振峰更尖锐、$V_{0m}$更接近$V_m$；$R$增大时，并联品质因子降低，峰值下降、频率选择性变差。两种情况的本质都是$R$影响$Q$值。[1]

    \item \textbf{力学实验要求"等振荡稳定后"记录数据，本实验为何不需要？}

    解：稳态建立时间约为$5\tau$。对RLC电路，$\tau = 2L/R = 2 \times 0.02/2500 = 16\,\mu\mathrm{s}$，故$5\tau \approx 80\,\mu\mathrm{s}$。这意味着从信号加入到达到稳态只需约$80\,\mu\mathrm{s}$，远快于示波器的刷新周期（通常$>10\,\mathrm{ms}$），因此人眼观察到的波形已是稳态结果。

    而力学弹簧振子的$Q$值可达数十甚至上百，$\tau$为数秒到数十秒量级。初始瞬态需要较长时间才能衰减，因此必须"等振荡稳定后"才能记录有效的稳态数据。[2]
\end{enumerate}

% ============ 八、参考文献 ============
\section*{参考文献}

\noindent [1] 清华大学物理系. 基础物理实验讲义：RLC电路[Z]. 北京：清华大学, 2024.

\noindent [2] Karni S. Applied Circuit Analysis[M]. New York: John Wiley \& Sons Inc., 1988: 220--233.

% ============ 九、附录 ============
\section*{附录}

\subsection*{原始数据}

原始实验数据如图~\ref{fig:raw}所示。

\begin{figure}[H]
\centering
\includegraphics[width=0.9\textwidth]{Fig17.png}
\caption{实验原始数据记录}
\label{fig:raw}
\end{figure}

\end{document}
